%Each project has its own unique flow, would recommend you write a sample flow (like a quick table of contents with sections \& sub-sections) of the Literature Review chapter and ask for feedback from your supervisor and implement it.\\
%But general idea:\\
%\begin{itemize}
%    \item Purpose?? It is to provide a comprehensive and critical overview of existing literature on a topic.\\
%    \item How?? By summarising , evaluating, and synthesising the relevant research and articles.\\
%    \item Why?? To identify gaps, trends, contradictions, and key findings. BUT also to justify your methods.
%\end{itemize}
%To start off the section you can provide a small intro and summary on what the section discusses\ldots

\section{Introduction}
%Introduce the purpose and scope of the literature review.

This chapter gives a critical overview of existing research on prompt injection
 (PI) and retrieval-augmented generation (RAG) poisoning. It summarises and 
 evaluates the relevant studies before drawing their findings together. 
 The aim is to identify gaps and trends in how adversaries manipulate large 
 language models (LLMs), and in how current defence mechanisms attempt to 
 secure these systems.

\section{Background}
%Provide a comprehensive background on the chosen topic (e.g., blockchain, AI).
%Explain the evolution and context of the topic.

According to \citet{correia_systematic_2026, greshake_not_2023} Large Language Models (LLMs) 
have rapidly transitioned from standalone text generation interfaces into 
foundational cognitive engines integrated into mission-critical software, 
enterprise workflows, and autonomous decision pipelines. Despite their linguistic fluency and advanced problem-solving capabilities,
foundation models suffer from well-documented systemic constraints,
including temporal knowledge obsolescence, factual hallucinations, 
and an inability to access private or domain-specific corpora natively 
\citet{zou_poisonedrag_2025,zhang_benchmarking_2025}.
To address these limitations, modern artificial intelligence
architectures widely adopt Retrieval-Augmented Generation (RAG) 
\citet{lewis2021retrievalaugmentedgenerationknowledgeintensivenlp}.
By coupling generative models with neural information retrieval pipelines
over external knowledge bases, RAG grounds model generation in dynamic, 
verifiable external evidence.

However, incorporating external data and third-party tools directly into 
the generative execution loop creates severe vulnerabilities
\citet{correia_systematic_2026}.In classical computing architectures, 
security relies fundamentally on strict, hardware-enforced boundaries 
separating executable control instructions from untrusted application data
\citet{greshake_not_2023}. In contrast, generative language models process 
system instructions, user queries, retrieved context, and tool outputs 
within a single, unified natural language context window
\citet{liu_formalizing_2024}.
Because transformers cannot fundamentally distinguish the authoritative
semantic weight of a system prompt from ingested external text, adversaries 
can manipulate the system's execution flow without modifying underlying 
model weights \citet{liu_formalizing_2024}.
This vulnerability manifests across two converging threat vectors:
\begin{enumerate}
    \item 
    \textbf{Prompt Injection (PI):} The insertion of adversarial natural 
    language commands through direct user prompts or indirect external 
    data channels (e.g., websites, emails, shared documents) that hijack 
    the model's instruction-following capabilities, overriding developer 
    directives and forcing unauthorized actions 
    \citet{greshake_not_2023,liu_formalizing_2024}. 
    \item 
    \textbf{RAG Knowledge Poisoning:} The targeted corruption of external retrieval 
    databases, wherein an adversary injects maliciously crafted texts 
    engineered to be retrieved by dense neural retrievers, thereby 
    manipulating the generator model into producing false statements, 
    commercial biases, or denial-of-service refusals
    \citet{zou_poisonedrag_2025,zhang_benchmarking_2025}.

\end{enumerate}




\section{Problem Description}
Elaborate on the specific problem identified in the introduction.

\section{Sources \& Impacts}
Explore relevant sources and literature related to the problem.
Discuss the impact of the problem on the field or industry.

\section{Existing Research / Solutions}
Review existing research and solutions addressing the identified problem.
Identify strengths and weaknesses of current approaches.

\section{Comparative Analysis}
Compare different projects, algorithms, solutions, or research work to your project title. Think about using a framework to better structure the comparison.

\section{Critical Analysis}
Analyse the strengths and limitations of existing research/projects/solutions critically.

\section{Summary of Findings}
Summarise the key findings from the literature review.

\section{Proposed Solution}
Introduce the proposed solution or approach based on the literature reviewed and gaps identified.
